\documentclass[11pt]{article}
\usepackage[a4paper,margin=18mm]{geometry}
\usepackage{fontspec}
\setmainfont{Latin Modern Roman}
\usepackage{amsmath,amssymb,amsthm,amscd,graphicx,color}
\usepackage{xurl}
\usepackage[unicode,hidelinks]{hyperref}
\allowdisplaybreaks
\setlength\emergencystretch{3em}
\numberwithin{equation}{section}

\newtheorem{Theorem}{Theorem}[section]
\newtheorem{Lemma}[Theorem]{Lemma}
\newtheorem{Proposition}[Theorem]{Proposition}
 { \theoremstyle{definition}
\newtheorem{Definition}[Theorem]{Definition}
\newtheorem{Remark}[Theorem]{Remark} }


\begin{document}
\begin{center}\Large Identity component of holomorphic automorphisms of the rank-one Deligne-Hitchin moduli space for a compact curve of genus at least two\end{center}
\noindent\textbf{Stable ID:} 1886\par
\noindent\textbf{Authors:} Indranil Biswas; Sebastian Heller\par
\noindent\textbf{Original work:} On the Automorphisms of a Rank One Deligne-Hitchin Moduli Space\par
\noindent\textbf{Version:} Official SIGMA 13 (2017) article 072, DOI 10.3842/SIGMA.2017.072; journal PDF and native TeX acquired 2026-10-01, exact bytes pinned by SHA256\par
\noindent\textbf{Selected task scope:} X compact connected complex Riemann surface of genus g>=\allowbreak{}2, MDH(X) the rank-one degree-zero Deligne-Hitchin space glued from X and its conjugate by the stated Riemann-Hilbert map. Only Aut\_hol(MDH)\_0 =\allowbreak{} MdR(X) times (H0(X,KX) times H0(X,KX)) semidirect C\textasciicircum{}times, with multiplication exactly3.7 (opposite weights on the two form factors).\par
\noindent\textbf{Difficulty:} H2: The classification first forces arbitrary holomorphic automorphisms to preserve the twistor fibers using normal-bundle degree and interpolation by explicit sections. It then solves the global gluing and gauge-lattice constraints to identify the translation factors and the semidirect scaling action. The proof relies on the geometry of moduli spaces and twistor sections rather than an immediate homogeneous-space automorphism formula.\par
\noindent\textbf{Editorial boundary note (not author text):} The complete original Theorem 3.7 and every new section-lift, fiber-preservation and gauge-lattice argument are retained for the selected genus>=\allowbreak{}2 rank-one degree-zero Deligne–Hitchin space. Exact Riemann–Hilbert gluing, all section lifts/constant harmonic gauge differences, the normal-degree contradiction, integral-lattice residual automorphisms and existence/uniqueness of translation are supplied; the original opposite-weight group law is unchanged. Section 2.1 Hodge-space definitions, projection and full author construction are retained as uncounted supporting context. The general hyperkahler twistor normal/tangent-bundle computation is an explicitly cited standard prerequisite, with the source model and degree formula preserved. Full disconnected/cohomology-preserving automorphisms and Moishezon consequences are unselected.\par
\noindent\textbf{Source:} \url{https://sigma-journal.com/2017/072/}\quad DOI: \url{https://doi.org/10.3842/SIGMA.2017.072}\par
\noindent\textbf{License and attribution:} Copyright retained by the named authors. Published by SIGMA under CC BY 4.0: \url{https://creativecommons.org/licenses/by/4.0/}. Source: \url{https://sigma-journal.com/about.html}. Adaptation consists of standalone excerpt formatting, cover and numbering restoration; original author language and mathematical text are retained.\par
\medskip\hrule\medskip\noindent\textbf{Author text begins below.}\par
\setcounter{section}{1}
\setcounter{subsection}{0}
\setcounter{subsubsection}{0}
\setcounter{Theorem}{5}
\setcounter{equation}{0}
% Original source lines 116--256
\section{Automorphisms of the rank one Hodge moduli space}\label{auto-hodge}

\subsection{Connected component of the automorphism group}

Let $X$ be an irreducible smooth complex projective curve of genus $g$, with $g \geq 1$. The holomorphic cotangent bundle of $X$ will be denoted by $K_X$.
\begin{Definition} For any $\lambda \in \mathbb C$, a rank one $\lambda$-connection on $X$ is a pair of the form $(L, D)$, where $L$ is a holomorphic line bundle on~$X$ and
\begin{gather*}
D \colon \ L \longrightarrow L\otimes K_X
\end{gather*}
is a holomorphic dif\/ferential operator of order no more than one such that
\begin{gather*}
D(f_0s) = f_0D(s)+\lambda\cdot s\otimes (df_0)
\end{gather*}
for every locally def\/ined holomorphic section $s$ of $L$ and every locally def\/ined holomorphic function $f_0$ on $X$.
\end{Definition}

So, if $(L, D)$ is a $\lambda$-connection with $\lambda \not= 0$, then $D/\lambda$ is a holomorphic connection on $L$; a~$0$-connection $D$ on $L$ is an element of $H^0(X, K_X)$ because such a $D$ is ${\mathcal O}_X$-linear. Note that for a~$\lambda$-connection $(L, D)$, with $\lambda \not= 0$ we have $\text{degree}(L) = 0$ because $L$ admits a holomorphic connection. Also, note that any holomorphic connection on a Riemann surface is automatically f\/lat.

If $(L, D)$ is a $0$-connection we will always impose the condition that $\text{degree}(L) = 0$.

Let
\begin{gather*}
{\mathcal M}_{\rm Hod} = {\mathcal M}^X_{\rm Hod}(1)
\end{gather*}
be the Hodge moduli space of rank one $\lambda$-connections. So, the points of ${\mathcal M}_{\rm Hod}$ parametrize triples of the form $(\lambda, L, D)$, where $\lambda \in \mathbb C$, $L$ is a holomorphic line bundle on $X$ of degree zero and~$D$ is a $\lambda$-connection on~$L$.

Def\/ine an algebraic morphism
\begin{gather}\label{lf}
f \colon \ {\mathcal M}_{\rm Hod} \longrightarrow {\mathbb C} , \qquad (\lambda, L, D) \longmapsto \lambda .
\end{gather}
For any $\lambda \in \mathbb C$, let
\begin{gather*}%\label{fiber}
f^\lambda := f^{-1}(\lambda) \subset {\mathcal M}_{\rm Hod}
\end{gather*}
be the f\/iber over $\lambda$. By def\/inition, the f\/iber $f^0$ is the moduli space of Higgs line bundles of degree zero $\operatorname{Pic}^0(X)\times H^0(X, K_X)$ on $X$, and the f\/iber $f^1$ is the moduli space of rank one holomorphic connections on $X$. The latter space is also called the de Rham moduli space. For $\lambda \not= 0$, the f\/iber $f^\lambda$ is canonically identif\/ied with $f^1$ by the map $(L, D) \longmapsto (L, D/\lambda)$. We note that $f^1$ is biholomorphic to the Betti moduli space $\operatorname{Hom}(\pi_1(X), {\mathbb C}^*) = ({\mathbb C}^*)^{2g}$ by sending a holomorphic connection to its monodromy representation, and hence $f^1$ does not admit a~compact submanifold of positive dimension. Since $\operatorname{Pic}^0(X)\times H^0(X, K_X)$ has the projective variety $\operatorname{Pic}^0(X)$ of positive dimension as a~subvariety, it follows that~$f^0$ is not biholomorphic to~$f^1$.

The group of all algebraic automorphisms of the quasiprojective variety
${\mathcal M}_{\rm Hod}$ will be denoted by $\operatorname{Aut}({\mathcal M}_{\rm Hod})$. Let
\begin{gather*}
\operatorname{Aut}({\mathcal M}_{\rm Hod})_0 \subset \operatorname{Aut}({\mathcal M}_{\rm Hod})
\end{gather*}
be the connected component containing the identity automorphism and
 let
\begin{gather*}
{\mathbb V} := \operatorname{Morphisms}\big({\mathbb C}, H^0(X, K_X)\big)
\end{gather*}
be the inf\/inite dimensional complex vector space parametrizing all algebraic maps from $\mathbb C$ to the af\/f\/ine space $H^0(X, K_X)$. Then we have the following theorem.
\begin{Theorem}\label{thm1}
The group $\operatorname{Aut}({\mathcal M}_{\rm Hod})_0$ fits in a short exact sequence of groups
\begin{gather*}
0 \longrightarrow {\mathbb V} \stackrel{\iota}{\longrightarrow} \operatorname{Aut}({\mathcal M}_{\rm Hod})_0 \stackrel{h}{\longrightarrow}
\operatorname{Pic}^0(X)\times {\mathbb C}^* \longrightarrow 0 .
\end{gather*}
\end{Theorem}

\begin{proof} Let
\begin{gather*}
T \colon \ {\mathcal M}_{\rm Hod} \longrightarrow {\mathcal M}_{\rm Hod}
\end{gather*}
be any algebraic automorphism that lies in $\operatorname{Aut}({\mathcal M}_{\rm Hod})_0$. It is known that the f\/iber $f^1$ does not admit any nonconstant algebraic function \cite[Proposition~2.2]{BBS}. Since $f^\lambda$ is isomorphic to $f^1$ for every $\lambda \in {\mathbb C}^*$, it follows that $f^\lambda$ does not admit any nonconstant algebraic function if $\lambda \in {\mathbb C}^*$. Hence the composition
\begin{gather}\label{e1}
f^\lambda \stackrel{T\vert_{f^\lambda}}{\longrightarrow} {\mathcal M}_{\rm Hod} \stackrel{f}{\longrightarrow} \mathbb C
\end{gather}
is a constant function for all $\lambda \in {\mathbb C}^*$. Note that the point $f\circ T(f^{\lambda})$ lies in ${\mathbb C}^*$ because $f^0$ is not isomorphic to $f^\lambda$. Since this also holds for~$T^{-1}$, it follows that the composition in~\eqref{e1} is a~constant function for each $\lambda \in {\mathbb C}$. This implies that there is an algebraic automorphism
\begin{gather*}
\tau_0 \colon \ {\mathbb C} \longrightarrow \mathbb C
\end{gather*}
such that
\begin{gather*}
f\circ T = \tau_0 \circ f .
\end{gather*}
Since $f^0$ is not isomorphic to any $f^\lambda$, $\lambda \in {\mathbb C}^*$, it follows that $\tau_0(0) = 0$. Therefore, $\tau_0$ coincides with the multiplication of $\mathbb C$ by a f\/ixed nonzero number; let
\begin{gather}\label{e4}
\tau \in {\mathbb C}^*
\end{gather}
be such that $\tau_0(z) = \tau \cdot z$ for all $z \in \mathbb C$.

The group of all automorphisms of the variety $\operatorname{Pic}^0(X)$ will be denoted by $\operatorname{Aut}(\operatorname{Pic}^0(X))$ (since the variety $\operatorname{Pic}^0(X)$ is projective, any holomorphic automorphism of it is algebraic). The connected component of~$\operatorname{Aut}(\operatorname{Pic}^0(X))$ containing the identity automorphism will be denoted by~$\operatorname{Aut}(\operatorname{Pic}^0(X))_0$. The automorphisms of $\operatorname{Pic}^0(X)$ given by translations of the group $\operatorname{Pic}^0(X)$ lie in $\operatorname{Aut}(\operatorname{Pic}^0(X))_0$. In fact, this way $\operatorname{Pic}^0(X)$ gets identif\/ied with
$\operatorname{Aut}(\operatorname{Pic}^0(X))_0$; note that the Lie algebra $H^0(\operatorname{Pic}^0(X), T\operatorname{Pic}^0(X))$ of~$\operatorname{Aut}(\operatorname{Pic}^0(X))_0$ coincides with the Lie algebra of $\operatorname{Pic}^0(X)$ by evaluating sections of~$T\operatorname{Pic}^0(X)$ at the identity element.

For any $\lambda \in \mathbb C$, let $\phi^\lambda$ be the morphism def\/ined by
\begin{gather}\label{e7}
\phi^\lambda \colon \ f^\lambda \longrightarrow \operatorname{Pic}^0(X) ,\qquad (L, D) \longmapsto L .
\end{gather}
The f\/ibers of $\phi^\lambda$ are isomorphic to $H^0(X, K_X)$, because the space of all holomorphic connections on a line bundle $L \in \operatorname{Pic}^0(X)$ is an af\/f\/ine space for the vector space $H^0(X, K_X)$, and $H^0(X, K_X)$ is the space of all Higgs f\/ields on any line bundle. There is no nonconstant algebraic map from an af\/f\/ine space to an abelian variety. Hence, there is no nonconstant algebraic map from a f\/iber of $\phi^\lambda$ to $\operatorname{Pic}^0(X)$. Consequently, we get a map
\begin{gather*}
\Phi \colon \  {\mathbb C} \longrightarrow \operatorname{Aut}\big(\operatorname{Pic}^0(X)\big),
\end{gather*}
which is uniquely determined by the condition that the following diagram is commutative
\begin{gather*}
\begin{matrix}
f^\lambda & \stackrel{T\vert_{f^\lambda}}{\longrightarrow} & f^{\tau\lambda}\\
~ ~~ \Big\downarrow \phi^\lambda && ~ ~~~~ \Big\downarrow \phi^{\tau\lambda}\\
\operatorname{Pic}^0(X) & \stackrel{\Phi(\lambda)}{\longrightarrow} &\operatorname{Pic}^0(X)
\end{matrix}
\end{gather*}
for all $\lambda \in \mathbb C$, where $\tau$ is the complex number in \eqref{e4}$ $. Note that since $T \in \operatorname{Aut}({\mathcal M}_{\rm Hod})_0$, it follows that the image of $\Phi$ lies in $\operatorname{Aut}\big(\operatorname{Pic}^0(X)\big)_0 = \operatorname{Pic}^0(X) \subset \operatorname{Aut}\big(\operatorname{Pic}^0(X)\big)$. Thus, $\Phi$ is a~constant map, again because there is no nonconstant algebraic map from $\mathbb C$ to $\operatorname{Pic}^0(X)$. Let{\samepage
\begin{gather}\label{e5}
\widehat{\Phi} = \Phi({\mathbb C}) \in \operatorname{Pic}^0(X)
\end{gather}
be the image of the map $\Phi$.}

Let
\begin{gather*}
h \colon \ \operatorname{Aut}({\mathcal M}_{\rm Hod})_0 \longrightarrow \operatorname{Pic}^0(X)\times {\mathbb C}^*
\end{gather*}
be the map def\/ined by $T \longmapsto \big(\widehat{\Phi}, \tau\big)$. It is straight-forward to check that $h$ is a group homomorphism.

We will now show that $h$ is surjective. For this, f\/irst note that the multiplicative action of~${\mathbb C}^*$ on $\mathbb C$ has a natural lift to an action of ${\mathbb C}^*$ on ${\mathcal M}_{\rm Hod}$. Indeed, any $c \in {\mathbb C}^*$ acts on ${\mathcal M}_{\rm Hod}$ as
\begin{gather*}
(\lambda, L, D) \longmapsto (c\cdot\lambda, L, c\cdot D) .
\end{gather*}
For all these automorphisms of ${\mathcal M}_{\rm Hod}$ given by the action of ${\mathbb C}^*$, the corresponding elements of~$\operatorname{Pic}^0(X)$ (see \eqref{e5}) coincide with the identity element. Therefore, to prove that~$h$ is surjective, it suf\/f\/ices to show that the composition of $h$ with the projection $\operatorname{Pic}^0(X)\times {\mathbb C}^*  \longrightarrow \operatorname{Pic}^0(X)$ is surjective. Take any $L_0 \in \operatorname{Pic}^0(X)$ and f\/ix a holomorphic connection~$D_0$ on~$L_0$. Def\/ine
\begin{gather*}
\beta \colon \  {\mathcal M}_{\rm Hod} \longrightarrow {\mathcal M}_{\rm Hod} ,\qquad
(\lambda, L, D) \longmapsto (\lambda, L\otimes L_0, (D\otimes \text{Id}_{L_0}) +(\lambda\cdot \text{Id}_L\otimes D_0)) .
\end{gather*}
It is straight-forward to check that $\beta \in \operatorname{Aut}({\mathcal M}_{\rm Hod})$; its inverse is the corresponding map for $(L_0^*, D_0^*)$. Since the moduli space of rank one connections on $X$ is connected, it follows that $\beta \in \operatorname{Aut}({\mathcal M}_{\rm Hod})_0$; note that $\beta$ for the trivial connection $({\mathcal O}_X, d)$ is the identity map of ${\mathcal M}_{\rm Hod}$. The element in $\operatorname{Pic}^0(X)$ corresponding to $\beta$ (see~\eqref{e5}) is clearly $L_0$. Therefore, the composition of $h$ with the projection $\operatorname{Pic}^0(X)\times {\mathbb C}^*  \longrightarrow \operatorname{Pic}^0(X)$ and hence $h$ itself are surjective.

Next, we need to def\/ine ${\mathbb V} \stackrel{\iota}{\longrightarrow} \operatorname{Aut}({\mathcal M}_{\rm Hod})_0$. To do so, consider for $v \in \mathbb V$ the automorphism
\begin{gather}\label{e6}
\iota_v \colon \ {\mathcal M}_{\rm Hod} \longrightarrow {\mathcal M}_{\rm Hod} ,\qquad (\lambda, L, D) \longmapsto (\lambda, L, D+v(\lambda)) .
\end{gather}
Clearly, we have $\iota_v \in \operatorname{Aut}({\mathcal M}_{\rm Hod})_0$, and also $\iota_v \in \operatorname{kernel}(h)$. Therefore, there is an injective homomorphism
\begin{gather*}
\iota \colon \ {\mathbb V} \longrightarrow \operatorname{kernel}(h) \subset \operatorname{Aut}({\mathcal M}_{\rm Hod})_0 ,\qquad v \longmapsto \iota_v .
\end{gather*}
We will prove that $\operatorname{image}(\iota) = \operatorname{kernel}(h)$.

In order to prove that $\operatorname{image}(\iota) = \operatorname{kernel}(h)$, take any $T \in \operatorname{Aut}({\mathcal M}_{\rm Hod})_0$ such that
\begin{gather*}T \in \operatorname{kernel}(h) .\end{gather*} We have $T(f^\lambda) = f^\lambda$ for all
$\lambda \in {\mathbb C}^*$ because the constant in \eqref{e4} for $T$ is~$1$. Since the element in \eqref{e5} corresponding to $T$ is the trivial line bundle, we get a morphism
\begin{gather*}
\delta_\lambda \colon \  f^\lambda \longrightarrow H^0(X, K_X) ,\qquad z \longmapsto T(z) - z ;
\end{gather*}
note that since the f\/ibers of $\phi^\lambda$ (constructed in \eqref{e7}) are af\/f\/ine spaces for $H^0(X, K_X)$, and $\phi^\lambda(T(z)) = \phi^\lambda(z)$ for all $z \in f^\lambda$, we have $T(z) - z \in H^0(X, K_X)$. As there are no nonconstant algebraic functions on $f^\lambda$ \cite[Proposition~2.2]{BBS}, it follows that the above function $\delta_\lambda$ is a constant one.

All automorphism of the af\/f\/ine space $H^0(X, K_X)$ are of the form $u \longmapsto A(u)+u_0$, where $A \in \text{GL}(H^0(X, K_X))$ and $u_0 \in H^0(X, K_X)$. Considering the restrictions of $T$ to open subsets of the form $h^{-1}(V\times U)$, where $U$ is an analytic neighborhood of $0 \in \mathbb C$ and $V$ is an analytic open subset of $\operatorname{Pic}^0(X)$, it now follows that the above map
\begin{gather*}\lambda \longmapsto \operatorname{image}(\delta_\lambda) \in H^0(X, K_X)\end{gather*}
extends across $0 \in \mathbb C$ as a holomorphic map from $\mathbb C$ to $H^0(X, K_X)$. In other words, there is a~unique element $v \in {\mathbb V}$ such that $v(\lambda) = \operatorname{image} (\delta_\lambda)$ for all $\lambda \in {\mathbb C}^*$. Clearly, we have $\iota_v = T$, where~$\iota_v$ is constructed in~\eqref{e6}. This implies that $\operatorname{image}(\iota) = \operatorname{kernel}(h)$. This completes the proof of the theorem.
\end{proof}

\setcounter{section}{2}
\setcounter{subsection}{2}
\setcounter{subsubsection}{0}
\setcounter{Theorem}{2}
\setcounter{equation}{7}
% Original source lines 273--275
\begin{gather}\label{phi}
\phi \colon \ {\mathcal M}_{\rm Hod} \longrightarrow \operatorname{Pic}^0(X) ,\qquad (\lambda, L, D) \longmapsto L ;
\end{gather}

\setcounter{section}{2}
\setcounter{subsection}{2}
\setcounter{subsubsection}{0}
\setcounter{Theorem}{3}
\setcounter{equation}{14}
% Original source lines 348--532
\section[Holomorphic automorphisms of the Deligne--Hitchin moduli space]{Holomorphic automorphisms of the Deligne--Hitchin\\ moduli space}

\subsection{Connected component of the holomorphic automorphism group}\label{aut0mdh}

The smooth locus of the moduli space of f\/lat connections on a compact Riemann surface is equipped with a natural hyper-K\"ahler structure; see~\cite{H} for the case of ${\rm SL}(2,\mathbb C)$-connections, and \cite{GX} for a detailed treatment of the case of f\/lat line bundles.

A hyper-K\"ahler manifold is a Riemannian manifold $(M, g)$ which is K\"ahler for three complex structures $I$, $J$ and $K$ satisfying the quaternionic relations
\begin{gather*}IJ = -JI = K .\end{gather*} The interplay between the Riemannian, symplectic and complex geometric aspects of hyper-K\"ahler manifolds makes them fascinating mathematical objects. A hyper-K\"ahler mani\-fold~$M$ admits a twistor space which contains all of the information about $M$ in a~complex geometric fashion: The compatible complex structures on $M$ are parametrized by $\mathbb CP^1 = S^2 = \big\{(x_1, x_2, x_3) \in {\mathbb R}^3 \,|\, x^2_1+x^2_2+x^2_3 = 1\big\}$ via
\begin{gather*}(x_1, x_2, x_3) \longmapsto I_{(x,_1,x_2,x_3)} = x_1 I+x_2 J+x_3 K ,\end{gather*}
where $I$, $J$, $K$ are the three original complex structures on~$M$. The twistor space $\mathcal Z$ is a~complex manifold with a holomorphic surjective submersion to $\mathbb CP^1$ such that the f\/iber over every \smash{$z \in {\mathbb C}P^1$} is identif\/ied with~$(M, I_z)$. The twistor lines are the ``constant'' sections of the above f\/ibration ${\mathcal Z} = M\times\mathbb CP^1 \longrightarrow {\mathbb C} P^1$. These sections are holomorphic and the normal bundle of a~twistor line is isomorphic to ${\mathcal O}_{\mathbb CP^1}(1)^{\oplus d} \longrightarrow \mathbb CP^1$, where $d = \frac{1}{2}\dim_{\mathbb R}M$. This normal bundle is equipped with additional structures which provide all of the information needed to reconstruct the hyper-K\"ahler structure on~$M$; see~\cite{HKLR} for details.

For a complex reductive group $G_{\mathbb C}$, the hyper-K\"ahler structure on the moduli space of f\/lat $G_{\mathbb C}$-connections on a compact Riemann surface $X$ is given by non-abelian Hodge theory. The complex structure $J$ is induced by the complex Lie group $G_{\mathbb C}$ through the identif\/ication of the moduli space with the Betti moduli space $\operatorname{Hom}(\pi_1(X), G_{\mathbb C})/\!\!/G_{\mathbb C}$. Via the solutions of Hitchin's self-duality equation, the moduli space of stable $G_{\mathbb C}$-Higgs bundles, i.e., the Dolbeault moduli space, gets identif\/ied with the moduli space of irreducible f\/lat connections, yielding the complex structure~$I$. The reverse map ${\mathcal M}_{dR}  \longrightarrow \mathcal M_{\rm Dol}$ arises from Donaldson's twisted harmonic maps construction \cite{Do} for the case of $G_{\mathbb C}={\rm SL }(2,\mathbb C)$, see~\cite{Co} for the general case. It was f\/irst shown by Hitchin in the case of ${\rm SL}(2,\mathbb C)$ that~$I$,~$J$ and $K := IJ$ give rise to a~hyper-K\"ahler structure for the natural Riemannian metric. In the abelian case of $G_{\mathbb C} = \mathbb C^*$ the theory simplif\/ies considerably boiling down to the classical abelian Hodge theory~\cite{GX}.

It was f\/irst noticed by Deligne that the twistor space of the moduli space of f\/lat connections on a Riemann surface $X$ has a convenient description by gluing the Hodge moduli spaces for~$X$ and~$\overline X$ via the Riemann--Hilbert isomorphism; see~\cite{Si95,Si08} for details. We recall below this construction for the rank one case (the case of our interests).

Take a point $(\lambda, L, D) \in \mathcal M_{\rm Hod}(X) = \mathcal M^X_{\rm Hod}(1)$ with $\lambda \neq 0$ and consider the f\/lat connec\-tion~$D/\lambda$ on $L$. Let
\begin{gather*}\rho \colon \ \pi_1(X)  \longrightarrow H^1(X, {\mathbb Z}) \longrightarrow \mathbb C^*\end{gather*}
be the monodromy representation for~$D/\lambda$. After identifying $H^1(X, {\mathbb Z})$ with $H^1(\overline{X}, {\mathbb Z})$ by the identity map of the underlying $C^\infty$ manifolds, the homomorphism~$\rho$ gives a~holomorphic rank one $1$-connection $(1, E, \widetilde D)$ on~$\overline{X}$. Consider the morphism~$f$ in~\eqref{lf}. Let
\begin{gather*}
f_{\overline X} \colon \ {\mathcal M}_{\rm Hod}(\overline{X}) := {\mathcal M}^{\overline X}_{\rm Hod}(1)  \longrightarrow {\mathbb C}
\end{gather*}
be the morphism obtained by substituting $\overline X$ in place of~$X$ in \eqref{lf}. Def\/ine a holomorphic map
\begin{gather}\label{hodge-gluing}
\varphi \colon \  f^{-1}({\mathbb C}^*)  \longrightarrow f^{-1}_{\overline X}({\mathbb C}^*) ,\qquad (\lambda, L, D) \longmapsto \big(\lambda^{-1}, E, \widetilde{D}/\lambda\big)
\end{gather}
of the restrictions of the f\/iber bundles ${\mathcal M}_{\rm Hod}(X)$ and ${\mathcal M}_{\rm Hod}(\overline{X})$ to ${\mathbb C}^* \subset \mathbb C$. This $\varphi$ is clearly a~biholomorphism.

\begin{Definition}\label{de1} The rank one Deligne--Hitchin moduli space \begin{gather*}{\mathcal M}_{\rm DH} = {\mathcal M}_{\rm DH}(X)\end{gather*} is
\begin{gather*}{\mathcal M}_{\rm DH}(X) = \mathcal M_{\rm Hod}(X)\cup_\varphi\mathcal M_{\rm Hod}(\overline{X}).\end{gather*}
\end{Definition}

We note that the map $f\vert_{f^{-1}({\mathbb C}^*)}$ coincides with the composition
\begin{gather*}
f^{-1}({\mathbb C}^*) \stackrel{\varphi}{\longrightarrow} f^{-1}_{\overline X}({\mathbb C}^*) \stackrel{f_{\overline X}}{\longrightarrow}
 {\mathbb C}^* \stackrel{z\mapsto z^{-1}}{\longrightarrow} {\mathbb C}^* .
\end{gather*}
Therefore, $f$ and the composition
\begin{gather*}
\mathcal M_{\rm Hod}(\overline{X}) \stackrel{f_{\overline X}}{\longrightarrow}  {\mathbb C} \stackrel{z\mapsto z^{-1}}{\longrightarrow}
{\mathbb C}P^1\setminus\{0\}
\end{gather*}
patch together to produce a morphism
\begin{gather}\label{mapf}
f \colon \ {\mathcal M}_{\rm DH} \longrightarrow \mathbb CP^1 .
\end{gather}
For any $\lambda \in \mathbb CP^1$, the f\/iber $f^{-1}(\lambda)$ will be denoted by $f^\lambda$. The f\/iber $f^{\infty}$ is the moduli space of Higgs line bundles of degree zero on~$\overline X$. This $f$ is the twistor f\/ibration mentioned earlier.

\begin{Definition} The {\it degree} of a smooth map $i \colon \Sigma \longrightarrow {\mathcal M}_{\rm DH}(X)$ from a compact oriented surface $\Sigma$ is the degree of the composition $f\circ i \colon \Sigma \longrightarrow {\mathbb C}P^1$, where $f$ is the projection in~\eqref{mapf}.
\end{Definition}

\begin{Lemma}\label{lemdeg} Let $\Sigma$ be a compact Riemann surface of genus $g_\Sigma$ and $i \colon \Sigma \longrightarrow {\mathcal M}_{\rm DH}(X)$ be a~holomorphic immersion. Then the degree of the corresponding holomorphic normal bundle \begin{gather*}{\mathcal N} = (i^*T{\mathcal
M}_{\rm DH})/T\Sigma\end{gather*} is
\begin{gather*}
\deg ({\mathcal N}) = (2 g_X +2)\deg (i)-\deg (T\Sigma) = (2 g_X +2)\deg (i) +2g_\Sigma-2 ,\end{gather*}
where $g_X$ is the genus of $X$.
\end{Lemma}

\begin{proof} As in the computation in \cite[pp.~555--556]{HKLR} it can be shown that the tangent bundle of the complex manifold ${\mathcal M}_{\rm DH}$ is holomorphically isomorphic to the pull-back
\begin{gather*}f^*\big({\mathcal O}_{{\mathbb C}P^1}(2)\oplus \big({\mathcal O}_{{\mathbb C}P^1}(1)^{\oplus 2g_X}\big)\big) ,\end{gather*}
where the f\/irst summand is the tangent bundle of $\mathbb CP^1$ and the direct sum ${\mathcal O}_{{\mathbb C}P^1}(1)^{\oplus 2g_X}$ corresponds to the hyper-K\"ahler structure on $\mathcal M_{dR}$. Now the lemma follows immediately from the fact that $\text{degree}(i^*f^*{\mathcal O}_{{\mathbb
C}P^1}(1)) = \text{degree}(i)$.
\end{proof}

We will compute $\operatorname{Aut}(\mathcal M_{\rm DH})_0$ by proving a series of lemmas.
\begin{Lemma}\label{sections}
Let $\alpha, \beta, \omega, \eta \in H^0(X, K_X)$ be holomorphic $1$-forms. There exists a holomorphic section $s \colon {\mathbb C}P^1 \longrightarrow {\mathcal M}_{\rm DH}$ defined by
\begin{gather*}s(\lambda) = [\lambda, \overline{\partial}+\overline{\omega}+\lambda\overline{\eta}, \lambda(\partial+\alpha)+\beta] \in \mathcal M_{\rm Hod}(X)
 \subset {\mathcal M}_{\rm DH}\end{gather*}
for $\lambda \in {\mathbb C} \subset \mathbb CP^1$. Conversely, every section of $f \colon {\mathcal M}_{\rm DH} \longrightarrow {\mathbb C}P^1 $ is of this form.
\end{Lemma}

\begin{proof}
The f\/irst part of the lemma is easily verif\/ied. In order to prove the converse direction we start with a lift of the family of holomorphic structures on ${\mathbb C} \subset {\mathbb C}P^1$: As $\mathbb C$ is simply connected, there exists a~map
\begin{gather*}
\lambda \longmapsto \overline{\omega}_1(\lambda) \in \overline{H^0(X, K_X)} , \qquad \lambda \in \mathbb C\end{gather*}
such that \begin{gather*}\phi\circ s(\lambda) = [\overline{\partial}+\overline{\omega}_1(\lambda)]\end{gather*}
for all $\lambda \in {\mathbb C}$, where~$\phi$ is the projection in \eqref{phi} and $[-]$ denotes gauge equivalence class. Let $\beta \in H^0(X, K_X)$ be the (well-def\/ined) Higgs f\/ield of $s(0)$. A lift of $s$ is then given by
\begin{gather*}\widehat{s} (\lambda) = (\lambda,\overline{\partial}+\overline{\omega}_1(\lambda), \beta+\lambda(\partial+\alpha_1(\lambda))\end{gather*}
for some holomorphic map \begin{gather*}\alpha_1 \colon \ \mathbb C \longrightarrow \Omega^{(1,0)}(X) .\end{gather*} The image of $\alpha_1$ is contained in $H^0(X, K_X) \subset \Omega^{(1,0)}(X)$ because the $\lambda$-connections are integrable. Thus, $\lambda \longmapsto s(\lambda)$, $\lambda \in \mathbb C^*$, produces the following family of f\/lat connections on~$X$:
\begin{gather*}\widehat{\nabla}^\lambda = d+\lambda^{-1}\beta+\overline{\omega}_1(\lambda)+ \alpha_1(\lambda) .\end{gather*}

We can repeat this over $\mathbb CP^1\setminus\{0\}$ for the Riemann surface $\overline X$. Indeed, there exists a lift $\widetilde s$ of the form
\begin{gather*}\widetilde{s}(\mu) = (\mu,\partial +\alpha_2(\mu), \mu(\overline{\partial}+ \overline{\omega}_2(\mu))+\overline{\eta})\end{gather*}
with $\mu = \frac{1}{\lambda}$. The corresponding f\/lat connections are denoted by $\widetilde{\nabla}^\lambda$. By the gluing condition in Def\/inition \ref{de1} there exists a~gauge transformation $g(\lambda) \colon X \longrightarrow \mathbb C^*$ for every $\lambda \in \mathbb C^*$ such that \begin{gather*}\widehat{\nabla}^\lambda\cdot g(\lambda) = \widetilde{\nabla}^\lambda ;\end{gather*}
this gauge transformation is unique up to a constant scalar. Since the connection 1-forms of $\widehat{\nabla}^\lambda$ and $\widetilde{\nabla}^\lambda$ are both harmonic 1-forms, their dif\/ference
\begin{gather*}\widehat{\nabla}^\lambda-\widetilde{\nabla}^\lambda = \chi(\lambda)\end{gather*}
is a lattice point
\begin{gather}\label{Gamma}
\chi(\lambda) \in \Lambda := \left\{\!\chi \in \Omega^1(X, \mathbb C) \,|\, d\chi = d*\chi = 0;
\int_\gamma\chi\in2\pi\sqrt{-1}\mathbb Z \text{ for all closed curves } \gamma\!\right\}\!.\!\!\!\!\end{gather}
As the connection 1-forms depend continuously on~$\lambda$, it follows that $\chi$ is independent of $\lambda$. By construction, the connection 1-form of $\lambda \longmapsto \widetilde{\nabla}^\lambda$ has a f\/irst order pole at $\lambda = \infty$ whose residue is~$\overline{\eta}$.
This implies that $\overline{\omega}_1(\lambda)$ must be a polynomial of degree at most one and that $\alpha_1(\lambda)$ is a~constant~$\alpha$.
\end{proof}

\begin{Lemma}\label{tau-lemma} Let $T \colon {\mathcal M}_{\rm DH} \longrightarrow {\mathcal M}_{\rm DH}$ be a~holomorphic automorphism. Then $T$ maps fibers of~$f$ to fibers of~$f$. Moreover, it covers the automorphism $\lambda \longmapsto \tau\lambda$ or $\lambda \longmapsto \frac{\tau}{\lambda}$ of ${\mathbb C}P^1$ for some $\tau \in \mathbb C^*$; if the latter case occurs, then the Jacobian of $X$ is isomorphic to the Jacobian of~$\overline X$.
\end{Lemma}

\begin{proof} We f\/irst prove that two points in the same f\/iber $p_1, p_2 \in f^{\lambda_0} = f^{-1}(\lambda_0)$ are mapped by~$T$ into a~single f\/iber~$f^{\lambda_1}$. To this end, we claim that any two points~$p_1$,~$p_2$ in dif\/ferent f\/ibers can be joined by a holomorphic section ${\mathbb C}P^1 \longrightarrow {\mathcal M}_{\rm DH}$ of the map~$f$. This can be proved quite easily by linear interpolation with sections of the form given in Lemma~\ref{sections}.

If two points $x_1$, $x_2$ of one f\/iber of $f$ are mapped by~$T$ into two dif\/ferent f\/ibers of $f$, we consider a holomorphic section~$s$ of~$f$ passing through~$T(x_1)$ and~$T(x_2)$. Then $\lambda \longmapsto T^{-1}(s(\lambda))$ is a~holomorphic immersion of~${\mathbb C}P^1$ to ${\mathcal M}_{\rm DH}$, and the degree of its normal bundle is the same as the degree of the normal bundle of the section~$s$. From Lemma~\ref{lemdeg} we obtain
\begin{gather*}(2 g_X+2)\deg (i)-2=2g_X,\end{gather*}
which yields $\deg (i)=1$. On the other hand, its degree is at least two because it passes through two points $x_1$, $x_2$ lying in one f\/iber. In view of this contradiction we conclude that~$T$ maps f\/ibers of~$f$ to f\/ibers of~$f$.

Therefore, there is a unique holomorphic map
\begin{gather*}T'\colon \ {\mathbb C}P^1 \longrightarrow {\mathbb C}P^1\end{gather*}
such that $f\circ T = T'\circ f$. As in the proof of Theorem~\ref{thm1}, $T$ must satisfy $T(\{0,\infty\}) = \{0,\infty\}$, and hence $T'$ is of the form $\lambda \longmapsto \tau\lambda$ or $\lambda \longmapsto \frac{\tau}{\lambda}$ for some $\tau \in \mathbb C^*$. If $T(0) = \infty$ the rank one Higgs moduli spaces for $X$ and $\overline X$ are holomorphically isomorphic. Since there is no nonconstant holomorphic map from an abelian variety
to an af\/f\/ine space, any biholomorphism between the rank one Higgs moduli spaces for $X$ and $\overline X$ produces a~biholomorphism between~$\operatorname{Pic}^0(X)$ and~$\operatorname{Pic}^0(\overline{X})$.
\end{proof}

Let $T \colon {\mathcal M}_{\rm DH} \longrightarrow {\mathcal M}_{\rm DH}$ be a holomorphic automorphism that maps each f\/iber of $f$ to itself, meaning $T(f^\lambda) = f^\lambda$ for all $\lambda \in {\mathbb C}P^1$. Consider the projection
\begin{gather}\label{pi2}
\pi_2 \colon \ f^0  = \operatorname{Pic}^0(X)\times H^0(X, K_X) \longrightarrow H^0(X, K_X) .
\end{gather}
Then, the automorphism $T$ restricted to $f^0$ maps f\/ibers of $\pi_2$ to f\/ibers of $\pi_2$ since $\operatorname{Pic}^0(X)$ does not have any nonconstant holomorphic map to $H^0(X, K_X)$. A~corresponding statement is true for the f\/iber $f^\infty$. Thus, the assumptions in the following
lemma are natural.

\begin{Lemma}\label{almost-id} Let $T \colon {\mathcal M}_{\rm DH} \longrightarrow {\mathcal M}_{\rm DH}$ be a~holomorphic automorphism such that
$T(f^\lambda) = f^\lambda$ for all $\lambda \in \mathbb CP^1$. Assume that $T$ restricted to $\operatorname{Pic}^0(X)\times\{0\} \subset f^0$ is the
identity map, and $T$ restricted to $\operatorname{Pic}^0(\overline{X})\times\{0\} \subset f^\infty$ is the identity map. Then, there exists an integral harmonic $1$-form $\gamma \in \Lambda = H^1(X, \mathbb Z) \subset \operatorname{Harm}(X; \mathbb C)$ such that $T$ is given by
\begin{gather*}T((\lambda, E, D)) = (\lambda, E, D+\lambda\gamma') ,\end{gather*}
where $\gamma = \gamma'+\gamma''$ is the decomposition into $(1,0)$ and $(0,1)$ components.
\end{Lemma}

\begin{proof}
Consider the ``constant'' section of $f$
\begin{gather*}s(\lambda) = [\lambda, \overline{\partial}, \lambda \partial],\end{gather*}
where $d = \partial+\overline{\partial}$ is the decomposition into types. From Lemma~\ref{sections} and the assumption that~$T$ restricted to $\operatorname{Pic}^0(X)\times\{0\} \subset f^0$ is the identity map it follows that $T\circ s$ must be of the form
\begin{gather}\label{specialT}
T\circ s(\lambda) = (\lambda, \overline{\partial}+\gamma_1'', \lambda (\partial+\gamma_2'))
\end{gather}
for some $\gamma_1, \gamma_2 \in \Lambda$. By applying the gauge $\exp\big({-}\int\gamma_1\big)$, the identity in \eqref{specialT} is equivalent to
\begin{gather*}T\circ s(\lambda) = (\lambda, \overline{\partial}, \lambda (\partial+\gamma'))\end{gather*}
 for $\gamma = \gamma_2-\gamma_1$.

Since $T$ is continuous it follows that for every constant section of~$f$ def\/ined by $s_D(\lambda) = (\lambda, L, \lambda D)$, where $D$ is a~holomorphic connection on~$L$, the equality
\begin{gather*}T\circ s_D(\lambda) = (\lambda, L, \lambda (D+\gamma'))\end{gather*}
holds. It should be mentioned that this is equivalent to
\begin{gather*}T\circ s(\lambda) = (\lambda, L\otimes L(\overline{\partial}-\gamma''), \lambda D) .\end{gather*}
Note that for any $\lambda \in {\mathbb C}\setminus\{0\}$ and for every point $p = (\lambda, E, D) \in f^\lambda$, there exists a constant section
which goes through $p$. This completes the proof.
\end{proof}

Let \begin{gather*}s \colon \ {\mathbb C}P^1 \longrightarrow {\mathcal M}_{\rm DH} ,\qquad
\lambda \longmapsto [\lambda, \overline{\partial} +\overline{\omega}+ \lambda\overline{\eta}, \lambda(\partial+\alpha)+\beta]
\end{gather*}
be a section of $f$, where $\alpha, \beta, \eta, \omega \in H^0(X, K_X)$; def\/ine
\begin{gather}\label{Ts}
T_s \colon \ {\mathcal M}_{\rm DH} \longrightarrow {\mathcal M}_{\rm DH} ,\qquad [\lambda, E, D] \longmapsto [\lambda, E\otimes
L(\overline{\partial}+\overline{\omega}+\lambda\overline{\eta}), D\otimes(\lambda(\partial+\alpha)+\beta)] ,
\end{gather}
where the tensor product of two $\lambda$-connections~$D$,~$\widetilde{D}$ on two holomorphic line bundles~$E$,~$\widetilde E$ respectively is the
$\lambda$-connection
\begin{gather*}D\otimes\widetilde{D}(e\otimes\widetilde{e}) = D(e)\otimes \widetilde{e}+e\otimes \widetilde{D}(\widetilde {e})\end{gather*}
on $E\otimes\widetilde{E}$. This $T_s$ is clearly a holomorphic automorphism of ${\mathcal M}_{\rm DH}$.

Let
\begin{gather*}
\operatorname{Aut}({\mathcal M}_{\rm DH})_0 \subset \operatorname{Aut}({\mathcal M}_{\rm DH})
\end{gather*}
be the connected component, containing the identity element, of the group of holomorphic automorphisms of ${\mathcal M}_{\rm DH}$. Using Lemmas~\ref{tau-lemma} and~\ref{almost-id} we can give the following description of $\operatorname{Aut}({\mathcal M}_{\rm DH})_0$.

\begin{Theorem}\label{thmDHa}
There is an isomorphism
\begin{gather*}\operatorname{Aut}({\mathcal M}_{\rm DH})_0 =
\mathcal M_{dR}(X)\times \big(\big(H^0(X, K_X)\times\overline{H^0(X, K_X)}\big)\rtimes {\mathbb C}^*\big)\end{gather*}
with the group structure of the right-hand side given by
\begin{gather}\label{gs37}
\big(\nabla^1, \alpha_1, \overline{\eta}_1, \tau_1\big)\cdot \big(\nabla^2, \alpha_2,\overline{\eta}_2, \tau_2\big) = \big(\nabla^1\otimes\nabla^2, \alpha_1+\tau_1\alpha_2, \overline{\eta}_1+\tfrac{1}{\tau_1}\overline{\eta}_2, \tau_1\tau_2\big) .
\end{gather}
The group $\mathcal M_{dR}(X)$ acts via the automorphisms $T_s$ in~\eqref{Ts}, $H^0(X, K_X)\times\overline{H^0(X, K_X)}$ acts using addition of forms to connections and Higgs bundles, and ${\mathbb C}^*$ acts via multiplication.
\end{Theorem}

\begin{proof} Using Lemma \ref{tau-lemma} and the lift of the $\mathbb C^*$ action on $\mathbb C P^1$ we can restrict to automorphisms $T\in \operatorname{Aut}(\mathcal M_{\rm DH})_0$ such that
\begin{gather*}T\big(f^\lambda\big) = f^\lambda\end{gather*} for all $\lambda\in\mathbb CP^1$. We claim that there exists a section $s \colon {\mathbb C}P^1 \longrightarrow {\mathcal M}_{\rm DH}$ of $f$ such that $T_s\circ T$ is of the form given in Lemma~\ref{almost-id}, where $T_s$ is constructed in~\eqref{Ts}.

Recall that $T$ acts f\/iberwise on the f\/ibers of~$\pi_2$ in~\eqref{pi2}. An automorphism of $\operatorname{Pic}^0(X)$ which is homotopic to the identity map is a translation by an element of $\operatorname{Pic}^0(X)$. Hence, there exists $\alpha, \omega \in H^0(X, K_X)$ such that
\begin{gather*}T(0, L, 0) = (0, L\otimes L(\overline{\partial}+\overline{\omega}), \alpha)\end{gather*}
for all $(0, L, 0) \in \pi_2^{-1}(0) \subset f^0$. The same argument yields that there are $\beta, \eta \in H^0(X, K_X)$ such that the automorphism $T_s$, for the section $s$ def\/ined by
\begin{gather*}s(\lambda) = (\lambda, \overline{\partial} -\overline{\omega}-\lambda\overline{\eta}, \lambda(\partial-\beta)-\alpha) ,\end{gather*}
has the desired properties.

Moreover, by applying the arguments in the proof of Lemma~\ref{almost-id}, we see that a~section~$s$ of~$f$ such that $T_s\circ T$ is of the form given in Lemma~\ref{almost-id} is unique up to tensoring with a constant section of the form $\lambda \longmapsto (\lambda, \overline{\partial}, \lambda(\partial+\gamma'))$ for some $\gamma \in H^1(X, \mathbb Z) \subset \operatorname{Harm}(X, \mathbb C)$.

Altogether, we obtain that any element in $\operatorname{Aut}({\mathcal M}_{\rm DH})_0$ is of the form $(\nabla, \alpha, \overline{\eta}, \tau)\in\mathcal M_{dR}(X)\times \big(\big(H^0(X, K_X)\times\overline{H^0(X, K_X)}\big)\rtimes {\mathbb C}^*\big)$ for the actions of the components $\mathcal M_{dR}(X)$, $H^0(X, K_X)$, $\overline{H^0(X, K_X)}$ and ${\mathbb C}^*$ on ${\mathcal M}_{\rm DH}$ as described above. The group structure of $\operatorname{Aut}({\mathcal M}_{\rm DH})_0$ is easily verif\/ied to be as in \eqref{gs37}.
\end{proof}
\begin{thebibliography}{99}
\bibitem[3]{BBS}
Baraglia D., Biswas I., Schaposnik L.P., Automorphisms of {$\mathbb{C}^*$}
 moduli spaces associated to a {R}iemann surface, \href{https://doi.org/10.3842/SIGMA.2016.007}{\textit{SIGMA}} \textbf{12}
 (2016), 007, 14~pages, \href{https://arxiv.org/abs/1508.06587}{arXiv:1508.06587}.

\bibitem[4]{Co}
Corlette K., Flat {$G$}-bundles with canonical metrics, \href{https://doi.org/10.4310/jdg/1214442469}{\textit{J.~Differential
 Geom.}} \textbf{28} (1988), 361--382.

\bibitem[5]{Do}
Donaldson S.K., Twisted harmonic maps and the self-duality equations,
 \href{https://doi.org/10.1112/plms/s3-55.1.127}{\textit{Proc. London Math. Soc.}} \textbf{55} (1987), 127--131.

\bibitem[7]{GX}
Goldman W.M., Xia E.Z., Rank one {H}iggs bundles and representations of
 fundamental groups of {R}iemann surfaces, \href{https://doi.org/10.1090/memo/0904}{\textit{Mem. Amer. Math. Soc.}}
 \textbf{193} (2008), viii+69~pages, \href{https://arxiv.org/abs/math.DG/0402429}{math.DG/0402429}.

\bibitem[8]{H}
Hitchin N.J., The self-duality equations on a {R}iemann surface, \href{https://doi.org/10.1112/plms/s3-55.1.59}{\textit{Proc.
 London Math. Soc.}} \textbf{55} (1987), 59--126.

\bibitem[9]{HKLR}
Hitchin N.J., Karlhede A., Lindstr\"om U., Ro\v{c}ek M., Hyperk\"ahler metrics
 and supersymmetry, \href{https://doi.org/10.1007/BF01214418}{\textit{Comm. Math. Phys.}} \textbf{108} (1987), 535--589.

\bibitem[11]{Si95}
Simpson C., The {H}odge f\/iltration on nonabelian cohomology, in Algebraic
 Geometry~-- {S}anta {C}ruz 1995, \href{https://doi.org/10.1090/pspum/062.2/1492538}{\textit{Proc. Sympos. Pure Math.}}, Vol.~62,
 Amer. Math. Soc., Providence, RI, 1997, 217--281, \href{https://arxiv.org/abs/alg-geom/9604005}{alg-geom/9604005}.

\bibitem[12]{Si08}
Simpson C., A weight two phenomenon for the moduli of rank one local systems on
 open varieties, in From {H}odge Theory to Integrability and {TQFT}
 tt*-Geometry, \href{https://doi.org/10.1090/pspum/078/2483751}{\textit{Proc. Sympos. Pure Math.}}, Vol.~78, Amer. Math. Soc.,
 Providence, RI, 2008, 175--214, \href{https://arxiv.org/abs/0710.2800}{arXiv:0710.2800}.

\end{thebibliography}
\end{document}
